\chapter{Derivations from Statistical Mechanics}
\label{app:statmech}

These are the derivations underlying Chapter~\ref{ch:background}. They are
placed here so that the body carries the argument and the appendix carries
the work; nothing in them is omitted.

\section{Hamiltonian mechanics and phase space}\label{sec:hamiltonian}

Consider a mechanical system with $n$ degrees of freedom, described by
generalised coordinates $q = (q_1, \dots, q_n)$ and conjugate momenta
$p = (p_1, \dots, p_n)$. Its dynamics is generated by a single scalar
function, the \emph{Hamiltonian} $H(q, p)$, which for the systems
considered here has the mechanical form
\begin{equation}
  H(q, p) \;=\; \underbrace{\sum_{i=1}^{n} \frac{p_i^2}{2m_i}}_{\text{kinetic}}
  \;+\; \underbrace{U(q_1, \dots, q_n)}_{\text{potential}},
  \label{eq:hamiltonian}
\end{equation}
and equals the total energy of the system. The state of the system at any
instant is the pair $(q, p)$, a point in the $2n$-dimensional \emph{phase
space} $\Gamma \subseteq \R^{2n}$. The evolution is given by
\emph{Hamilton's equations},
\begin{equation}
  \dot{q}_i \;=\; \frac{\partial H}{\partial p_i},
  \qquad
  \dot{p}_i \;=\; -\,\frac{\partial H}{\partial q_i},
  \qquad i = 1, \dots, n.
  \label{eq:hamilton}
\end{equation}
Two structural properties of \eqref{eq:hamilton} are the pillars on which
statistical mechanics rests.

\paragraph{Energy conservation.} Along any trajectory,
\begin{equation}
  \frac{\dd H}{\dd t}
  = \sum_{i=1}^n \left(
      \frac{\partial H}{\partial q_i}\,\dot q_i
      + \frac{\partial H}{\partial p_i}\,\dot p_i
    \right)
  = \sum_{i=1}^n \left(
      \frac{\partial H}{\partial q_i}\frac{\partial H}{\partial p_i}
      - \frac{\partial H}{\partial p_i}\frac{\partial H}{\partial q_i}
    \right)
  = 0 .
  \label{eq:energy-conservation}
\end{equation}
The motion is therefore confined to the \emph{energy shell}
$\Gamma_E = \{(q,p) : H(q,p) = E\}$, a $(2n-1)$-dimensional hypersurface.

\paragraph{Incompressibility of the phase flow.} The phase-space velocity
field
\begin{equation}
  v(q, p) \;=\; (\dot q, \dot p)
  \;=\; \left( \frac{\partial H}{\partial p},\, -\frac{\partial H}{\partial q} \right)
  \label{eq:phase-velocity}
\end{equation}
has identically vanishing divergence:
\begin{equation}
  \nabla \cdot v
  \;=\; \sum_{i=1}^n \left(
      \frac{\partial \dot q_i}{\partial q_i}
      + \frac{\partial \dot p_i}{\partial p_i}
    \right)
  \;=\; \sum_{i=1}^n \left(
      \frac{\partial^2 H}{\partial q_i \partial p_i}
      - \frac{\partial^2 H}{\partial p_i \partial q_i}
    \right)
  \;=\; 0 .
  \label{eq:incompressible}
\end{equation}
Hamiltonian flow behaves like an incompressible fluid in phase space: it
can stretch and fold a region, but never change its volume. This identity
is the seed of the probabilistic construction that follows.

\section{Liouville's theorem}\label{sec:liouville}

A single trajectory is useless for macroscopic physics: for a gas,
$n \sim 10^{23}$ and neither the initial condition nor the solution of
\eqref{eq:hamilton} is accessible. The move that founds statistical
mechanics is to describe our \emph{ignorance} of the microstate by a
probability density $\rho(q, p, t)$ on phase space:
$\rho(q,p,t)\,\dd q\,\dd p$ is the probability that the system is in the
infinitesimal cell around $(q,p)$ at time $t$. Liouville's theorem states
how this density evolves.

\begin{theorem}[Liouville]\label{thm:liouville}
Under the Hamiltonian flow \eqref{eq:hamilton}, the phase-space density
obeys
\begin{equation}
  \frac{\partial \rho}{\partial t}
  \;=\; -\,\{\rho, H\}
  \;=\; \sum_{i=1}^n \left(
    \frac{\partial H}{\partial q_i}\frac{\partial \rho}{\partial p_i}
    - \frac{\partial H}{\partial p_i}\frac{\partial \rho}{\partial q_i}
  \right),
  \label{eq:liouville}
\end{equation}
where $\{\cdot,\cdot\}$ is the Poisson bracket. Equivalently, the density
is constant along trajectories: $\dd \rho / \dd t = 0$.
\end{theorem}

\begin{toolbox}{Proof of Liouville's theorem, step by step}
The density is locally conserved (probability mass is neither created nor
destroyed, only transported by the flow), so it satisfies the continuity
equation familiar from fluid dynamics:
\begin{equation*}
  \frac{\partial \rho}{\partial t} + \nabla \cdot (\rho\, v) = 0,
  \qquad v = (\dot q, \dot p).
\end{equation*}
Expand the divergence of the product using
$\nabla \cdot (\rho v) = v \cdot \nabla \rho + \rho\, (\nabla \cdot v)$:
\begin{equation*}
  \frac{\partial \rho}{\partial t}
  + \underbrace{\sum_{i=1}^n \left(
      \dot q_i \frac{\partial \rho}{\partial q_i}
      + \dot p_i \frac{\partial \rho}{\partial p_i}\right)}_{v \cdot \nabla \rho}
  + \rho \underbrace{\sum_{i=1}^n \left(
      \frac{\partial \dot q_i}{\partial q_i}
      + \frac{\partial \dot p_i}{\partial p_i}\right)}_{\nabla \cdot v}
  = 0 .
\end{equation*}
By the incompressibility identity \eqref{eq:incompressible} the last term
vanishes. Substituting Hamilton's equations
$\dot q_i = \partial H/\partial p_i$, $\dot p_i = -\partial H/\partial q_i$
into the middle term:
\begin{equation*}
  \frac{\partial \rho}{\partial t}
  + \sum_{i=1}^n \left(
      \frac{\partial H}{\partial p_i}\frac{\partial \rho}{\partial q_i}
      - \frac{\partial H}{\partial q_i}\frac{\partial \rho}{\partial p_i}
    \right) = 0
  \quad\Longleftrightarrow\quad
  \frac{\partial \rho}{\partial t} = -\{\rho, H\}.
\end{equation*}
Finally, the total (convective) derivative along a trajectory is
\begin{equation*}
  \frac{\dd \rho}{\dd t}
  = \frac{\partial \rho}{\partial t} + v \cdot \nabla \rho
  = \frac{\partial \rho}{\partial t} + \{\rho, H\}
  = 0 ,
\end{equation*}
which is the second form of the statement: an observer riding a
phase-space trajectory sees a constant local density. \hfill$\square$
\end{toolbox}

Liouville's theorem has an immediate consequence: \emph{any density that
is a function of conserved quantities only is stationary}. If
$\rho(q,p) = f\big(H(q,p)\big)$ for some function $f$, then
$\{\rho, H\} = f'(H)\,\{H, H\} = 0$, so $\partial \rho / \partial t = 0$.
Equilibrium statistical mechanics is the study of which $f$ nature
chooses.

\section{From dynamics to ensembles}\label{sec:ensembles}

An \emph{ensemble} is a probability distribution over microstates that is
stationary under the dynamics and compatible with the macroscopic
constraints we impose (fixed energy, fixed temperature, and so on). Two
ensembles matter for this thesis.

\subsection{The microcanonical ensemble}

For a perfectly isolated system, the energy is fixed at $E$ and
Liouville's theorem suggests the simplest choice compatible with it: a
uniform density on the energy shell,
\begin{equation}
  \rho_{\mathrm{mc}}(q, p)
  \;=\; \frac{1}{\Omega(E)}\,
        \delta\big(H(q,p) - E\big),
  \qquad
  \Omega(E) = \int_\Gamma \delta\big(H(q,p) - E\big)\,\dd q\,\dd p ,
  \label{eq:microcanonical}
\end{equation}
where $\Omega(E)$ counts the microstates at energy $E$. The postulate
that all accessible microstates are equally likely (the \emph{equal a
priori probability} postulate) is motivated, not proved, by ergodicity:
for ergodic dynamics, time averages along a single trajectory converge to
averages over the shell, so the uniform measure is the one a long
experiment effectively samples. Boltzmann's entropy attaches a number to
this count,
\begin{equation}
  S(E) \;=\; k_B \log \Omega(E),
  \label{eq:boltzmann-entropy}
\end{equation}
with $k_B$ the Boltzmann constant. Entropy is the logarithm of the number
of microscopic ways a macroscopic state can be realised.

\subsection{The canonical ensemble and the Boltzmann distribution}
\label{sec:canonical}

Isolated systems are an idealisation; the physically ubiquitous situation
is a system exchanging energy with a much larger environment (a
\emph{heat bath}) at temperature $T$. The stationary distribution of the
system alone is then no longer uniform: it is the
\emph{Boltzmann--Gibbs distribution},
\begin{equation}
  \boxed{\;
  \rho_\beta(x) \;=\; \frac{1}{Z(\beta)}\, e^{-\beta H(x)},
  \qquad
  Z(\beta) = \int e^{-\beta H(x)}\,\dd x,
  \qquad
  \beta = \frac{1}{k_B T},
  \;}
  \label{eq:boltzmann}
\end{equation}
where $x = (q,p)$ collects the microstate and $Z(\beta)$, the
\emph{partition function}, normalises the density. Equation
\eqref{eq:boltzmann} can be derived in two independent ways, and both
derivations are instructive enough to be carried out in full: the first
(physical) traces the exponential to the entropy of the bath; the second
(information-theoretic) characterises it as the least-committal
distribution compatible with a fixed mean energy, a perspective due to
\citet{jaynes1957information} that returns, almost verbatim, in modern
machine learning.

\begin{toolbox}{The Boltzmann distribution from the heat bath}
Let the total system (system $S$ together with bath $B$) be isolated with
total energy $E_{\mathrm{tot}}$, and let the coupling be weak, so
$H_{\mathrm{tot}} = H_S + H_B$. Apply the microcanonical postulate to the
total system: the probability that $S$ is in a \emph{specific} microstate
$x$ with energy $H_S(x) = \epsilon$ is proportional to the number of bath
microstates carrying the remaining energy,
\begin{equation*}
  \Prob(x) \;\propto\; \Omega_B\big(E_{\mathrm{tot}} - \epsilon\big)
  \;=\; \exp\!\left[ \tfrac{1}{k_B} S_B\big(E_{\mathrm{tot}} - \epsilon\big) \right],
\end{equation*}
using the definition \eqref{eq:boltzmann-entropy} of the bath entropy.
Since the bath is huge, $\epsilon \ll E_{\mathrm{tot}}$, expand $S_B$ to
first order:
\begin{equation*}
  S_B(E_{\mathrm{tot}} - \epsilon)
  = S_B(E_{\mathrm{tot}})
  - \epsilon \,\frac{\partial S_B}{\partial E}\Big|_{E_{\mathrm{tot}}}
  + O\!\big(\epsilon^2 / E_{\mathrm{tot}}\big).
\end{equation*}
The derivative of the bath entropy with respect to its energy
\emph{defines} the bath temperature,
$\partial S_B / \partial E = 1/T$. Substituting,
\begin{equation*}
  \Prob(x) \;\propto\;
  \underbrace{e^{S_B(E_{\mathrm{tot}})/k_B}}_{\text{constant in } x}
  \cdot\, e^{-\epsilon/(k_B T)}
  \;\propto\; e^{-\beta H_S(x)} .
\end{equation*}
The quadratic remainder is suppressed by the bath's heat capacity and
vanishes in the thermodynamic limit. Normalising yields
\eqref{eq:boltzmann}. Note what happened: the exponential form is nothing
but the first-order Taylor expansion of the bath's entropy. The Boltzmann
factor is the price, in bath microstates, of lending energy $\epsilon$ to
the system. \hfill$\square$
\end{toolbox}

\begin{toolbox}[tb:maxent]{The Boltzmann distribution from maximum entropy}
Forget physics; ask, with \citet{jaynes1957information}: among all
densities $\rho$ with a given mean energy $\E_\rho[H] = U$, which one
assumes the least beyond the constraint? ``Least'' is measured by the
Gibbs--Shannon entropy $S[\rho] = -k_B \int \rho \log \rho$. Maximise
$S[\rho]$ subject to $\int \rho = 1$ and $\int \rho H = U$ with Lagrange
multipliers $\lambda_0, \lambda_1$:
\begin{equation*}
  \mathcal{L}[\rho]
  = -k_B \!\int\! \rho \log \rho \,\dd x
  \;-\; \lambda_0 \Big( \!\int\! \rho \,\dd x - 1 \Big)
  \;-\; \lambda_1 \Big( \!\int\! \rho H \,\dd x - U \Big).
\end{equation*}
Take the functional derivative and set it to zero pointwise:
\begin{equation*}
  \frac{\delta \mathcal{L}}{\delta \rho(x)}
  = -k_B\big(\log \rho(x) + 1\big) - \lambda_0 - \lambda_1 H(x) = 0
  \;\;\Longrightarrow\;\;
  \rho(x) = \exp\!\Big[-1 - \tfrac{\lambda_0}{k_B} - \tfrac{\lambda_1}{k_B} H(x)\Big].
\end{equation*}
The density is an exponential of the constrained quantity. This is
completely general: constraining $\E[\phi_j(x)]$ for features $\phi_j$
produces $\rho \propto \exp[-\sum_j \beta_j \phi_j(x)]$, the exponential
family. The multiplier $\lambda_0$ is fixed by normalisation and produces
$Z$; renaming $\lambda_1/k_B = \beta$ and identifying $\beta = 1/(k_B T)$
through thermodynamics recovers \eqref{eq:boltzmann} exactly.
Second-order check:
$\delta^2 \mathcal{L} / \delta \rho^2 = -k_B/\rho < 0$, so the stationary
point is the unique maximum (the entropy is strictly concave).
\hfill$\square$
\end{toolbox}

The second derivation deserves a remark. The Boltzmann distribution is
not a statement about molecules but about \emph{inference under
constraints}. Any time we posit a model of the form
$p(x) \propto e^{-E(x)}$, as we will for Markov random fields
(Chapter~\ref{ch:graphs}) and energy-based models
(Appendix~\ref{app:diffusion-deriv}), we are doing statistical mechanics, whether the
variables are spins, pixels, or the frames of a video.

\section{The partition function and free energy}\label{sec:free-energy}

The partition function $Z(\beta)$ in \eqref{eq:boltzmann} looks like a
mere normaliser; it is in fact the central computational object of the
theory, and the reason approximate inference exists as a field. Its
logarithm defines the \emph{Helmholtz free energy},
\begin{equation}
  F(\beta) \;=\; -\frac{1}{\beta} \log Z(\beta).
  \label{eq:free-energy}
\end{equation}

\begin{toolbox}{Everything is a derivative of $\log Z$}
Differentiate $\log Z$ with respect to $\beta$, carrying the derivative
through the integral:
\begin{equation*}
  -\frac{\partial \log Z}{\partial \beta}
  = -\frac{1}{Z}\frac{\partial}{\partial \beta}\!\int\! e^{-\beta H}
  = \frac{\int H\, e^{-\beta H}}{\int e^{-\beta H}}
  = \E_{\rho_\beta}[H] \;=\; U ,
\end{equation*}
the mean energy. Differentiate once more:
\begin{align*}
  \frac{\partial^2 \log Z}{\partial \beta^2}
  &= \frac{\partial}{\partial \beta}
     \left( -\frac{\int H e^{-\beta H}}{Z} \right)
   = \frac{\int H^2 e^{-\beta H}}{Z}
     - \left( \frac{\int H e^{-\beta H}}{Z} \right)^{\!2} \\[2pt]
  &= \E[H^2] - \E[H]^2 \;=\; \operatorname{Var}(H) \;\geq\; 0 .
\end{align*}
So $\log Z$ is convex in $\beta$, its first derivative gives the mean
energy, its second the energy fluctuations; the latter equals
$k_B T^2 C_V$ with $C_V$ the heat capacity. This is a
\emph{fluctuation--dissipation} relation: how a system responds to
heating is determined by its spontaneous equilibrium fluctuations.

The same identity holds for any observable coupled linearly to a field. If the
Gibbs weight carries a factor $e^{h \cdot a}$, then
$Z(h) = \int \nu(a) e^{h\cdot a}\,\dd a$ is a cumulant generating function and
\begin{equation}
  \frac{\partial \langle a_k \rangle}{\partial h_j}
  = \frac{\partial^2 \log Z}{\partial h_j \partial h_k}
  = \operatorname{Cov}(a_k, a_j) ,
  \label{eq:static-response}
\end{equation}
so a susceptibility is a covariance. Susceptibilities defined this way are the
standard probe of phase transitions, and they are the instrument used by
\citet{bachtis2024cascade} to locate the transitions discussed in
Section~\ref{sec:statmech-learning}. Chapter~\ref{ch:ring} shows that the
response of a diffusion score to its own observations is exactly of the form
\eqref{eq:static-response}, with the observation playing the role of the field.
Equation~\eqref{eq:static-response} is a static identity within one fixed
measure, and is not the two-time dynamical fluctuation--dissipation theorem of
non-equilibrium statistical mechanics.
Finally, rewriting the Gibbs--Shannon entropy of $\rho_\beta$:
\begin{equation*}
  S[\rho_\beta]
  = -k_B \E\big[\log \rho_\beta\big]
  = -k_B \E\big[-\beta H - \log Z\big]
  = \frac{U}{T} + k_B \log Z ,
\end{equation*}
and multiplying by $-T$: $\; -TS = -U + F$, i.e.
\begin{equation*}
  F = U - TS .
\end{equation*}
The free energy is the exact trade-off between energy and entropy: at low
$T$ minimising $F$ means minimising energy (ordered states win), at high
$T$ it means maximising entropy (disordered states win). Phase
transitions are the non-analyticities of $F$ where the winner changes.
\hfill$\square$
\end{toolbox}

There is one more characterisation of $F$ worth recording, because it is
the statistical-mechanics form of the variational principle that
underlies both the training objective of diffusion models
(Chapter~\ref{ch:diffusion}) and the Bethe free energies of message
passing. For \emph{any} trial density $\mu$,
\begin{equation}
  \underbrace{\E_\mu[H] - \tfrac{1}{\beta} S[\mu]/k_B}_{\text{variational free energy } F[\mu]}
  \;=\; F \;+\; \tfrac{1}{\beta}\, \kl{\mu}{\rho_\beta}
  \;\;\geq\;\; F ,
  \label{eq:gibbs-variational}
\end{equation}
with equality iff $\mu = \rho_\beta$: the Boltzmann distribution is the
unique minimiser of the variational free energy, and the gap is exactly a
Kullback--Leibler divergence. Every variational method in machine
learning is a restriction of \eqref{eq:gibbs-variational} to a tractable
family of trial densities.


\section{The Hopfield dynamics descends the energy}\label{sec:hopfield-descent}

Moved here from Chapter~\ref{ch:statmech}, where the statement is used but the
two lines of algebra behind it are not the point.

Let $h_i = \sum_{j \neq i} J_{ij} \sigma_j$ be the local field on neuron $i$
and suppose the update \eqref{eq:hopfield-dynamics} flips it, so
$\sigma_i' = -\sigma_i$. Only the terms of $H$ containing $\sigma_i$ change,
hence
\begin{equation*}
  \Delta H
  \;=\; H(\sigma') - H(\sigma)
  \;=\; -\,(\sigma_i' - \sigma_i)\, h_i
  \;=\; -\,2\,\sigma_i'\, h_i .
\end{equation*}
The update sets $\sigma_i' = \operatorname{sign}(h_i)$, so
$\sigma_i' h_i = |h_i| \geq 0$ and $\Delta H \leq 0$, strictly whenever a flip
actually occurs. $H$ is bounded below on a finite configuration space, so the
dynamics terminates in a local minimum and \emph{every} initial condition falls
into an attractor.

The single-pattern case $P=1$, the Mattis model \citep{mattis1976solvable}, is
gauge-equivalent to the uniform ferromagnet: the relabelling
$\sigma_i \to \xi_i \sigma_i$ maps the stored pattern onto the all-up state.
It returns in Section~\ref{sec:cascade} as the exactly solvable teacher in a
modern analysis of how energy-based models learn.

\section{First- and second-order transitions}\label{sec:transition-order}

The order of a transition is the order of the derivative of the free energy
that first misbehaves.

At a \emph{first-order} transition the first derivative jumps: the order
parameter is discontinuous, latent heat is absorbed, the two phases coexist at
the transition, and correlation lengths stay finite. At a \emph{second-order}
transition the first derivative is continuous and the second diverges: the
order parameter grows continuously from zero, the susceptibility
\eqref{eq:static-response} and the correlation length diverge, and observables
acquire power laws.

The distinction is what makes the learning cascade of Section~\ref{sec:cascade}
detectable through susceptibilities rather than through jumps, and it is also
why the finite models of this thesis cannot exhibit either in the strict sense:
both notions require the thermodynamic limit, and a correlation length that
grows until it reaches the system size and then stops is a finite-size ceiling.

\section{The learning cascade in detail}\label{sec:cascade-detail}

Supporting Section~\ref{sec:cascade}. The chapter states that RBM training is
a cascade of second-order transitions and what that means for this thesis; the
mapping that establishes it, and the empirical evidence, are here.

gradient \eqref{eq:bm-learning}. The central idea is to track training
through the singular value decomposition of the weight matrix,
\begin{equation}
  W \;=\; \sum_{\gamma=1}^{r} w_\gamma\, u_\gamma \bar u_\gamma^\T ,
  \label{eq:rbm-svd}
\end{equation}
whose modes $(u_\gamma, \bar u_\gamma)$ act as candidate order
parameters: the magnetisation of the visible layer along mode $\gamma$,
$m_\gamma = u_\gamma \cdot v / \sqrt{N_v}$, measures whether the learned
Gibbs measure has developed structure in that direction. At
initialisation the weights are small and the model is a paramagnet: no
direction in configuration space is preferred. The claim of
\citet{bachtis2024cascade} is that training is an \emph{effective
cooling}: as the singular values grow, the model crosses a sequence of
genuine second-order phase transitions, one for each feature it learns,
each signalled by a diverging susceptibility and falling in the
mean-field (Curie--Weiss) universality class. The mechanism is visible in
a two-line calculation.

\begin{toolbox}{The retrieval transition of a one-mode RBM}
Following \citet{bachtis2024cascade}, take a single Gaussian hidden unit
of variance $1/N_v$ coupled to $N_v$ binary visibles through a weight
vector $w$: the energy is
$E(v, h) = -h\, w \cdot v + \tfrac{N_v}{2} h^2$. Integrating out the
Gaussian $h$ gives the visible marginal
\begin{equation*}
  p(v) \;\propto\; \exp\!\Big[ \tfrac{1}{2 N_v}\, (w \cdot v)^2 \Big] :
\end{equation*}
a Mattis/Hopfield measure (Section~\ref{sec:hopfield}) whose single
pattern is $w$ itself, made extensive by the $1/N_v$ variance scaling.
Write $w = \bar w\, \xi$ with $\xi \in \{-1,+1\}^{N_v}$ and let
$m = \tfrac{1}{N_v} \sum_i \xi_i v_i$ be the overlap; then
$p(v) \propto \exp[\tfrac{N_v \bar w^2}{2} m^2]$. Counting configurations
at fixed overlap ($\asymp \exp[N_v H(\tfrac{1+m}{2})]$ with $H$ the
binary entropy) gives the free-entropy density
\begin{equation*}
  \phi(m) \;=\; \frac{\bar w^2}{2}\, m^2 + H\!\Big(\frac{1+m}{2}\Big)
  \;=\; \log 2 + \frac{\bar w^2 - 1}{2}\, m^2 - \frac{m^4}{12} + O(m^6).
\end{equation*}
This is a Landau expansion: the quadratic coefficient changes sign at
$\bar w^2 = \norm{w}^2 / N_v = 1$. Below the threshold the maximiser is
$m = 0$ (paramagnet); above it two symmetric maximisers $\pm m^*$ appear
continuously (retrieval). The transition is second order with mean-field
exponents, and the susceptibility along $w$ diverges as
$\norm{w}^2/N_v \uparrow 1$. Growing a singular value during training is
therefore literally a quench through a Curie--Weiss critical point.
\hfill$\square$
\end{toolbox}

\paragraph{Why a cascade.} With data generated by a Hopfield teacher
whose two patterns are correlated combinations
$\xi^{1,2} = \eta_1 \pm \eta_2$ of orthogonal directions, the data
covariance is rank two with unequal strengths, and the early-time
gradient flow of the student's SVD amplitudes along $\eta_1$ and $\eta_2$
is exponential with \emph{two different rates}. The stronger direction
(the centre of mass of the clusters) reaches its critical magnitude
first: the learned measure splits into two lumps along $\pm\eta_1$. The
weaker direction keeps growing at its own rate and crosses criticality
later, splitting each lump again: four modes, the full teacher, retrieved
in hierarchical order. Each crossing is a distinct second-order
transition with its own critical time, its own diverging susceptibility,
and hysteresis below it. Learning proceeds coarse-to-fine not by accident
of optimisation but because the critical times are ordered by the
spectrum of the data.

\paragraph{Evidence on real data.} On MNIST, CelebA, and human-genome
data, \citet{bachtis2024cascade} track the susceptibilities of individual
SVD modes during training of binary--binary RBMs and observe the same
signatures: successive susceptibility peaks that sharpen with system
size, finite-size-scaling collapses consistent with mean-field exponents
$(\gamma, \nu) = (1, \tfrac12)$ at upper critical dimension four, and
hysteresis, the operational fingerprints of genuine transitions rather
than smooth crossovers. This confirms and sharpens the earlier picture of
RBM learning as progressive alignment with the data's principal
components \citep{decelle2017spectral}, of feature emergence as
condensation \citep{tubiana2017emergence}, and of sequential mode
learning in linear networks \citep{saxe2014exact}; a broad review of the
EBM--physics interface is \citet{carbone2025hitchhiker}.

